\documentclass[11pt]{article}

\usepackage[T1]{fontenc}
\usepackage[utf8]{inputenc}
\usepackage{lmodern}
\usepackage[margin=1.04in]{geometry}
\usepackage{amsmath,amssymb,mathtools,amsthm}
\usepackage{aliascnt}
\usepackage{mathrsfs,microtype,enumitem}
\usepackage{xcolor,hyperref}
\hypersetup{colorlinks=true,linkcolor=blue!45!black,citecolor=green!35!black,urlcolor=blue!55!black,
  pdftitle={Paper IV Repair Note: Normalizations, Fixed Parts, and the Relative Period Torsor},
  pdfauthor={Christopher D. Long; Antoine-Auguste Le Blanc}}
\numberwithin{equation}{section}
\setlength{\emergencystretch}{3em}
\allowdisplaybreaks
\newtheorem{theorem}{Theorem}[section]
\newaliascnt{lemma}{theorem}
\newtheorem{lemma}[lemma]{Lemma}
\aliascntresetthe{lemma}
\newaliascnt{proposition}{theorem}
\newtheorem{proposition}[proposition]{Proposition}
\aliascntresetthe{proposition}
\newaliascnt{corollary}{theorem}
\newtheorem{corollary}[corollary]{Corollary}
\aliascntresetthe{corollary}
\theoremstyle{definition}
\newaliascnt{definition}{theorem}
\newtheorem{definition}[definition]{Definition}
\aliascntresetthe{definition}
\newaliascnt{externalinput}{theorem}
\newtheorem{externalinput}[externalinput]{External input}
\aliascntresetthe{externalinput}
\theoremstyle{remark}
\newaliascnt{remark}{theorem}
\newtheorem{remark}[remark]{Remark}
\aliascntresetthe{remark}
\usepackage[capitalize,noabbrev]{cleveref}
\crefname{externalinput}{External input}{External inputs}
\Crefname{externalinput}{External input}{External inputs}
\newcommand{\Q}{\mathbb Q}
\newcommand{\C}{\mathbb C}
\newcommand{\A}{\mathbb A}
\newcommand{\Gm}{\mathbb G_m}
\newcommand{\Ocal}{\mathcal O}
\newcommand{\Ecal}{\mathcal E}
\newcommand{\Mcal}{\mathcal M}
\newcommand{\Tcal}{\mathcal T}
\newcommand{\Dcal}{\mathcal D}
\newcommand{\Ncal}{\mathcal N}
\newcommand{\Ccal}{\mathcal C}
\newcommand{\Pcal}{\mathcal P}
\newcommand{\MN}{\mathsf{MN}}
\newcommand{\EM}{\mathsf{EM}_{\mathrm N}}
\newcommand{\Perv}{\mathsf{Perv}}
\newcommand{\Conn}{\mathsf{Conn}}
\newcommand{\VecCat}{\mathsf{Vec}}
\newcommand{\rat}{\operatorname{rat}}
\newcommand{\im}{\operatorname{im}}
\newcommand{\Hom}{\operatorname{Hom}}
\newcommand{\Aut}{\operatorname{Aut}}
\newcommand{\Isom}{\operatorname{Isom}}
\newcommand{\Spec}{\operatorname{Spec}}
\newcommand{\Frac}{\operatorname{Frac}}
\newcommand{\Span}{\operatorname{span}}
\newcommand{\id}{\operatorname{id}}
\newcommand{\dR}{\mathrm{dR}}
\newcommand{\dRB}{\mathrm{dRB}}
\newcommand{\B}{\mathrm B}
\newcommand{\lis}{\mathrm{lis}}
\newcommand{\mot}{\mathrm{mot}}
\newcommand{\diff}{\mathrm{diff}}
\newcommand{\fix}{\mathrm{fix}}
\newcommand{\PV}{\mathrm{PV}}
\newcommand{\pH}{{}^{p}\!H}
\newcommand{\one}{\mathbf1}
\newcommand{\DD}{\mathbb D}

\title{Paper IV Repair Note:\\
Normalizations, Fixed Parts, and the Relative Period Torsor}
\author{Christopher D. Long \and Antoine-Auguste Le Blanc\thanks{Le Blanc cryptographic author identifier: \texttt{b6048125\allowbreak{}30b3535b\allowbreak{}23d6dbb0\allowbreak{}f5abe32d\allowbreak{}1ed1a9de\allowbreak{}35fb4f7d\allowbreak{}04d62cff\allowbreak{}08ea19c8}}}
\date{September 2026}

\begin{document}
\maketitle

\begin{center}
\small
\texttt{galizur@gmail.com}
\end{center}

\begin{abstract}
This is a standalone set of replacement arguments for the pre-repair version
of Paper IV, \emph{Relative Exponential Fixed Parts and Functional Periods:
An Exponential Nori--Ayoub Theorem}. It fixes the exponential graph sign and
the passage from perverse to lisse tensor conventions; replaces two overly
general categorical arguments by affine-line adjunction; makes the unipotent
Fourier direct-image argument use actual exact sequences; and separates
algebraic and analytic fiber functors in the Galois argument. The relative
algebraic realization is constructed from the absolute Nori algebra and its
de Rham augmentation, with an explicit bar-resolution comparison and
six-operation verification. The derived kernel shift and the distinction
between topological and moderate nearby cycles are included. The last section
constructs the period comparison first over $\C(S)$, proves equivariance,
and then descends injectivity, differential simplicity, and the exact field
of constants. In particular, no assertion that the actual comparison matrix
is already Picard--Vessiot over the smaller period-constant field is used to
prove that assertion. The ordinary module-presentation,
absolute mixed-realization, Fourier, and $1$-motivic inputs are identified
separately from the new assembly and replacement proofs.
\end{abstract}

\paragraph{Frozen baseline and status.}
The repository path, frozen commit, and Paper IV source-blob identifier are,
respectively,
\begin{quote}\small\ttfamily
long-mathematics/exponential-integral-theorem\\
0bb151018dc73c9e48b847fff31fb82fd2cb9072\\
4cc45c6dd428e09cf5db940632106fc394d01393.
\end{quote}
The reference \path{baseline/paper-iv-pre-repair-2026-09-28} points to that
commit. The machine-readable baseline record contains the source and PDF
SHA-256 hashes. The frozen baseline and its 26-page PDF remain available at
that commit. The replacement arguments have received an assembled-proof
and integration audit and have been incorporated into the revised canonical
Paper IV on the repair branch. This note remains a separate development
companion; the dated closure record specifies the reviewed files and checks.
Neither that audit nor compilation is independent peer review or formal
verification.

\paragraph{Reading contract.}
The standard source theorems below are external inputs, not conclusions
of the desired period theorem. The previously outstanding relative
algebraic realization-and-comparison interface is now supplied by
\cref{thm:algebraic-interface,thm:algebraic-exponential,prop:relative-Betti-interface}.
Its construction is part of this repair note and was examined in the
assembled-proof audit recorded with the canonical integration. The final torsor theorem also has an abstract form:
once the fixed-part and rank-one hypotheses are supplied, its proof uses
no further motivic construction. This revision incorporates the kernel
and nearby-cycle corrections from the second-pass audit. Historical
reports retain their earlier verdicts; the ledger records the new proofs
without treating a successful build as a mathematical certification.

\tableofcontents

\section{Normalization conventions}\label{sec:normalization}

Put $k=\overline{\Q}\subset\C$. Let $S/k$ be smooth, connected, and
quasi-projective, of dimension $d$, and write $a:S\to\Spec k$ and
$K=k(S)$. Unless otherwise stated, $\lambda\in k^\times$. Geometric ground
field, coefficient field, de Rham field, and analytic constant field are not
silently identified: algebraic connections are defined over $k$, the
motivic categories are $k$-linear, and analytic comparison is taken after
extension to $\C$.

All six operations in this note use sheaf-compatible perverse
normalizations. Derived functors are derived even when no $R$ or $L$ is
written. A vector bundle with connection $M$ on a smooth $d$-fold corresponds
to a perverse-normalized object, while its unshifted de Rham complex has
terms in degrees $0,\ldots,d$.

\subsection{The exponential sign}

We distinguish the holonomic heart object from the derived kernel:
\begin{equation}\label{eq:exponential-sign}
 \Ecal^f=(\Ocal_S,d-df),\qquad
 \mathscr K_\lambda=(\Ocal_{\A^1_t},d-\lambda\,dt),\qquad
 \mathscr L_\lambda=\mathscr K_\lambda[-1].
\end{equation}
In left differential-operator notation,
$\partial_t\cdot p=p'-\lambda p$ on $\mathscr K_\lambda$.
A horizontal section is $e^{\lambda t}$; the dual exponential in an
integral pairing is $e^{-\lambda t}$. Virk defines the sheaf-compatible
derived kernel by requiring $\mathscr L_\lambda[1]$ to be this heart
object \cite[\S6.6 and Proposition~6.8]{Virk}.
The shift and the sign are separate conventions, both of which matter.
At the zero graph over a point, $i_0^*\mathscr L_\lambda=k$ in degree
zero, whereas $i_0^*\mathscr K_\lambda=k[1]$. This verifies the unit and
would detect an omitted kernel shift.

Let $\mathsf Q_S$ be the exponential quotient functor, and let
$i_h:S\to S\times\A^1_t$ be the graph of $h\in\Ocal(S)$.

\begin{lemma}[Graph realization]\label{lem:graph}
The normalized graph object
\[
  \mathsf G_S(h):=\mathsf Q_S(i_{h,*}\one_S[d])
\]
realizes the connection $\Ecal^{\lambda h}$ at scale $\lambda$. In
particular, the graph potential realizing $\Ecal^f$ is $h=f/\lambda$.
\end{lemma}
\begin{proof}
Let $p:S\times\A^1\to S$ and $\mu:S\times\A^1\to\A^1$ be the
projections. The algebraic realization constructed in
\cref{thm:algebraic-exponential} tensors with
$\mu^*\mathscr L_\lambda$ and applies $p_!$. Since $p i_h=\id_S$ and
the graph is proper, projection formula reduces it to
$\one_S[d]\otimes h^*\mathscr L_\lambda$. Its connection is
$d-\lambda\,dh$. The $[-1]$ already in $\mathscr L_\lambda$ gives the
unshifted rank-one local object under pullback; the remaining $[d]$ puts
that connection in the perverse heart. There is no extra shift to remove
and no change of sign.
\end{proof}

For example, on $\A^1_s$ at scale one, the potential $s$ gives $d-ds$,
not $d+ds$. The two connections are not isomorphic over $k(s)$: a rational
gauge between them would have logarithmic derivative $2$ or $-2$, whereas
a rational logarithmic derivative tends to zero at infinity.

\subsection{The lisse tensor normalization}

Write $\star$ for the derived additive convolution over $S$, including the
ordinary derived tensor product in the base direction. Its derived unit is
$\one_S^{\exp}:=\mathsf Q_S(i_{0,*}\one_S)$. For perverse-lisse objects put
\begin{equation}\label{eq:lisse-tensor}
  N\otimes_{\lis}M=(N\star M)[-d],\qquad
  \one_{S,\lis}=\one_S^{\exp}[d].
\end{equation}
The associativity, symmetry, and unit constraints are transported from the
unshifted lisse normalization by $N\mapsto N[-d]$ and its inverse.
This transport is part of the definition; the symmetry is not an
unadjusted restriction of the derived Koszul symmetry.

Fix a direction at infinity for the absolute Betti functor. For the
pairing with $e^{-\lambda t}$ choose the ray with
$\theta=-\arg\lambda$, or transport an equivalent choice consistently.
In this note $\Psi_\infty$ denotes this chosen directional functor
$\Psi_{e^{i\theta}\infty}$, not a direction-free canonical fiber; see
\cite[Definition~3.5 and Theorem~3.6]{Snodgrass}.

\begin{proposition}[Normalized tensor and fiber]\label{prop:normalization}
On the lisse subcategory the operations in \eqref{eq:lisse-tensor} realize
ordinary tensor products of connections. For an algebraic point $s\in S(k)$,
the normalized absolute restriction is $i_s^*N[-d]$. Its absolute
rapid-decay Betti fiber, in cohomological variance, is
\begin{equation}\label{eq:betti-fiber}
  \omega_{\B,s}(N)
  =\Psi_\infty\bigl(\rat^{\exp}(i_s^*N[-d])\bigr).
\end{equation}
In any strictly full rigid subquotient-closed lisse tensor subcategory this
is an exact tensor fiber functor. In homological notation one must dualize
consistently, rather than identify homology with cohomology.
\end{proposition}
\begin{proof}
For two connections, the sheaf-compatible derived tensor of their
perverse-normalized objects is their usual tensor connection shifted by
$d$. Shifting back by $[-d]$ therefore gives the usual tensor connection
in the heart. Realization compatibility with derived convolution gives the
first assertion, and reflection of the heart by exact faithful realization
gives lisseness of the result. The same calculation for the unit gives
$\one_{S,\lis}$.

A constant perverse-lisse object is $a^*C[d]$, and
$i_s^*a^*C[d][-d]=C$. For a general lisse object, normalized restriction
is the vector-space fiber of its connection, hence is in degree zero.
Exact faithful realization reflects that concentration, and the usual
fiber functor on connections is exact. Tensor compatibility follows by
applying $i_s^*[-d]$ to \eqref{eq:lisse-tensor}.
Finally the absolute nearby-fiber functor is exact, faithful, and tensor,
with the rapid-decay comparison; these are the absolute convolution inputs
of Katz and Fres\'an--Jossen, in the form recalled by
\cite[Corollary~2.17 and Theorem~3.6]{Snodgrass}.
\end{proof}

\begin{remark}[Fixed-scale duality]\label{rem:duality}
If $\epsilon(s,t)=(s,-t)$, then
$\rho_\lambda\epsilon^*=\rho_{-\lambda}$. Virk's actual duality formula is
\[
  \DD\rho_\lambda=\rho_{-\lambda}\DD,
\]
not same-scale commutation \cite[Theorem~6.13]{Virk}. On the rigid lisse
subcategory the usual expression for normalized tensor dual is
$\epsilon^*\DD N(-d)$: the Tate twist removes the smooth-base orientation
twist. Evaluation and coevaluation are understood in the transported
normalization. In particular the tensor inverse of $\mathsf G_S(h)$ is
$\mathsf G_S(-h)$; this follows directly from graph addition and does not
rely on a bare Verdier-duality assertion. We do not infer that every
object of the entire perverse heart is rigid.
\end{remark}

\section{The coefficient system and the affine-line quotient}\label{sec:heart}

\begin{externalinput}[Ordinary Nori structure and absolute comparison]
\label{input:ordinary}
We use the rational perverse-Nori coefficient system of
Ivorra--Morel, Terenzi, and Tubach, with its enhanced six operations,
exact faithful perverse Betti realization, and conservative derived Betti
realization \cite{IvorraMorel,Terenzi,Tubach}. Two precise additional
forms of these established inputs will be used.

First, put $\mathsf{DM}(X)=\mathsf{DM}_{\mathrm{et}}(X,\Q)$ and
$\mathsf{DN}(X)=\operatorname{Ind}D^b(\MN_{\Q}(X))$. If
$\mathsf{Nor}^*\dashv\mathsf{Nor}_*$ is Nori realization and
\[
 \mathscr N_k=\mathsf{Nor}_*\one,\qquad
 \mathscr N_X=a_X^*\mathscr N_k,
\]
then $\mathscr N_X\simeq\mathsf{Nor}_{X,*}\one$ and
\begin{equation}\label{eq:nori-module-model}
 \mathsf{DN}(X)\simeq
 \operatorname{Mod}_{\mathscr N_X}(\mathsf{DM}(X)).
\end{equation}
Under this symmetric monoidal equivalence, $\mathsf{Nor}_X^*Q$ is the
free module $\mathscr N_X\otimes Q$. The equivalence is compatible
with pullback and the constructible six operations. These are
Tubach's Lemma~4.12, Corollary~4.14, and Proposition~4.16; the compatibility
is part of the construction in his \S4.2.

Second, at $\Spec k$ ordinary Nori motives have the exact tensor de Rham
fiber $d_k$ over $k$, the rational Betti fiber $b_k$, and their natural
complex comparison. Their derived versions compute the corresponding
geometric realization complexes, compatibly with tensor products and
comparison. This is the \emph{absolute} mixed-realization construction,
not a relative realization assumption; see the cellular construction and
mixed-realization factorization in
\cite[Theorems~7.4.17 and~7.6.5]{Harrer} and
\cite[\S\S6--7, especially Propositions~7.1 and~7.12]{ChoudhuryGallauer}.
We use their complex-level constructions and extend them by Ind.
The covariant convention here is the homological one in
Choudhury--Gallauer; when using Harrer's contravariant cohomological
construction, take its dual on rigid geometric motives. This fixes
variance before any comparison of complexes.
\end{externalinput}

The missing relative algebraic realization is constructed below. In
particular, it is no longer included in \cref{input:ordinary} as a blanket
compatibility assumption. The construction uses the \emph{absolute}
Nori algebra in \eqref{eq:nori-module-model}; a choice of separate
algebraic models for complex Hodge objects would not give this construction.

\subsection{Algebraic geometric realization before the Nori quotient}

Let $\mathscr D_k(X)$ be the Ind-completion of the bounded constructible
regular-holonomic algebraic differential-module category on $X/k$, in
sheaf-compatible normalization. For singular $X$, use modules supported
on a local smooth ambient variety and Kashiwara equivalence. Transfer
bimodules and Spencer complexes give the enhanced operations. On a smooth
$d$-fold the unshifted tensor unit is $\Ocal_X[-d]$, not the holonomic
heart object $\Ocal_X$.

\begin{lemma}[Geometric de Rham coefficient system]\label{lem:geometric-dR}
There is a colimit-preserving symmetric monoidal geometric realization
\[
 r_X:\mathsf{DM}(X)\longrightarrow\mathscr D_k(X)
\]
which commutes with the six operations on constructible objects. After
$k\hookrightarrow\C$ and regular Riemann--Hilbert, it is the complexified
Betti realization, as a morphism of coefficient systems.
\end{lemma}
\begin{proof}
Use algebraic differential operators and their transfer bimodules over
$k$, rather than analytifying and choosing a descent afterward. Their
standard localization, base-change, projection, and smooth-purity maps
make $\mathscr D_k$ a rational stable coefficient system. These are the
ordinary algebraic differential-module operations; a realization model
retaining the filtered module over $k$ together with its complex
comparison is explicitly described in Saito's \S1.8(ii), with the
operations developed in his \S\S2--6 \cite{SaitoMixed}. No identification
of all such mixed sheaves with Nori motives is asserted here.

For clarity, the motivic localization conditions can be checked directly.
Etale descent is descent of the algebraic modules and transfer complexes.
For the affine-line projection, the relative algebraic de Rham complex of
the constant object is $\Ocal_X[t]\xrightarrow{\partial_t}
\Ocal_X[t]$, with cohomology $\Ocal_X$ in degree zero; integration of
polynomials is allowed in characteristic zero. Projection formula gives
$\id\simeq R\pi_*\pi^*$. The projective-line trace identifies the
reduced projective-line object with an invertible shifted Tate line, giving
stability. The usual de Rham orientation gives the rational motivic
factorization. These checks yield $r_X$ by the universal coefficient-system
construction \cite[Theorem~7.14]{DrewGallauer}. One restricts to the
geometric constructible image to obtain six-operation compatibility,
using the constructible realization argument of
\cite[Theorem~4.5 and its proof]{Tubach}; initiality alone is not being
used to assert compatibility with every right adjoint.

The relative de Rham theorem over $\C$, the trace map, and the Kunneth
map identify this system with Betti realization on smooth geometric
objects and their structure maps. Etale descent, affine-line localization,
and Tate stabilization extend this as a natural symmetric monoidal
comparison of the enhanced systems. The comparisons of adjunction and
base-change maps are the ones induced by the transfer complexes. Thus the
comparison is a comparison of functors and maps, not just of cohomology
dimensions or individual isomorphism classes.
\end{proof}

\begin{lemma}[The absolute comparison in the needed form]
\label{lem:absolute-dR-comparison}
On the absolute enhanced categories there is a symmetric monoidal
identification
\begin{equation}\label{eq:absolute-compatible}
 d_k\mathsf{Nor}_k^*\simeq r_k,
\end{equation}
and its complexification agrees with the actual absolute de Rham--Betti
comparison. The identification is compatible with the counit
$\mathsf{Nor}^*\mathsf{Nor}_*\one\to\one$.
\end{lemma}
\begin{proof}
Here is the specific use of the absolute input. For a good pair, its
Nori cohomology realizes as the algebraic relative de Rham cohomology and
as relative Betti cohomology with their integral comparison. Morphisms of
pairs, boundary maps, products, and Tate maps are included in this diagram.
Apply the two realizations to the good-filtration complexes used to
construct absolute derived Nori realization. They give the geometric
realization complexes with the same maps: the augmentation to the
geometric complex is the filtration spectral-sequence augmentation when
each relative pair has cohomology in its designated degree.

Use the system of good filtrations and common refinements, together with
finite affine Cech covers, not one independently chosen filtration for
each variety. Its refinement maps and product maps are the complex-level
construction in \cite[\S7]{ChoudhuryGallauer} and
\cite[Chapters~4 and~7]{Harrer}. Applying the exact realization functors
therefore gives a natural comparison of complexes, with the tensor
coherences supplied by that construction. Localization and Tate
stabilization give \eqref{eq:absolute-compatible}, and Ind extends it
from constructibles. This is the usual absolute mixed-realization
factorization in \cref{input:ordinary}. In particular, we do not need the
more delicate comparison of separately chosen \emph{relative} cellular
filtrations discussed in Tubach's Remark~4.10.

The de Rham--Betti comparison is natural on absolute Nori motives; apply
that naturality to the counit, viewed as a morphism in $\mathsf{DN}(k)$.
This proves the last assertion. No assertion that its matrix entries are
algebraic is required.
\end{proof}

\subsection{The absolute augmentation and the relative module formula}

\begin{lemma}[Extension from an augmented algebra]\label{lem:module-extension}
Let $\mathcal A$ and $\mathcal B$ be presentable stable symmetric monoidal
categories, let $r:\mathcal A\to\mathcal B$ preserve colimits and tensor
products, and let $T$ be a commutative algebra of $\mathcal A$ with an
augmentation $\alpha:r(T)\to\one_{\mathcal B}$. Then
\begin{equation}\label{eq:abstract-module-formula}
 F(M)=\one_{\mathcal B}\otimes^{\mathbf L}_{r(T)}r(UM)
\end{equation}
defines a colimit-preserving symmetric monoidal functor
$\operatorname{Mod}_{T}(\mathcal A)\to\mathcal B$. It sends the free
module $T\otimes Q$ to $r(Q)$. Among colimit-preserving $r$-linear functors with the specified
identification on free modules, the functor and its compatible natural
transformations are determined by $r$ and the action $\alpha$.
\end{lemma}
\begin{proof}
Applying $r$ to the multiplication and action maps makes $r(UM)$ an
$r(T)$-module. Derived extension of scalars along $\alpha$ gives
\eqref{eq:abstract-module-formula}; associativity of relative tensor
products gives its monoidal and colimit properties. For a free module the
formula reduces to
$\one\otimes_{r(T)}(r(T)\otimes r(Q))\simeq r(Q)$.

For the determination assertion use the augmented simplicial free bar
resolution of a $T$-module. Its terms are free modules on
$T^{\otimes n}\otimes UM$, and its face maps are multiplication and the
module action, with degeneracies induced by the unit. The underlying
augmented simplicial object is split by the unit, so its geometric
realization is $M$. A colimit-preserving $r$-linear functor with the given
action on the unit sends this bar construction to the relative-tensor
bar construction in \eqref{eq:abstract-module-formula}. This proves the
assertion for all modules and all module morphisms, not just for the
objects that happen to be free. It also proves that a compatible
comparison of $r$ and $\alpha$ induces a coherent comparison of $F$.
\end{proof}

By the lax symmetric monoidal right adjoint, $\mathscr N_k$ is a
commutative algebra. Define an algebra map entirely over $k$ by
\begin{equation}\label{eq:deRham-augmentation}
 \begin{split}
 \alpha_k:r_k(\mathscr N_k)
 &\simeq d_k(\mathsf{Nor}_k^*\mathsf{Nor}_{k,*}\one)\\
 &\xrightarrow{\ d_k(\varepsilon)\ }d_k(\one)=k.
 \end{split}
\end{equation}
Here $\varepsilon$ is the adjunction counit. For $a_X:X\to\Spec k$
put $\alpha_X=a_X^*\alpha_k$ under
$r_X(\mathscr N_X)\simeq a_X^*r_k(\mathscr N_k)$.
The map $\alpha_k$ is a \emph{de Rham augmentation}. It is not evaluation
of a numerical period algebra at an alleged algebraic period point.

\begin{theorem}[Algebraic realization and comparison interface]
\label{thm:algebraic-interface}
For quasi-projective $X/k$, under \eqref{eq:nori-module-model} set
\begin{equation}\label{eq:relative-algebraic-realization}
 R_{\dR,X}(M)=
 \one_{\mathscr D_k(X)}\otimes^{\mathbf L}_{r_X(\mathscr N_X)}
 r_X(UM),
\end{equation}
using $\alpha_X$. On constructibles this is an algebraic realization over
$k$ with the following properties.
\begin{enumerate}[label=\textup{(\roman*)},leftmargin=2.5em]
\item It is symmetric monoidal, compatible with the six operations and their
adjunction, projection, and base-change maps, and
$R_{\dR,X}\mathsf{Nor}_X^*=r_X$.
\item After $k\hookrightarrow\C$, regular Riemann--Hilbert identifies it
naturally with the complexified Betti realization. Equivalently, its
complex algebraic differential module is the underlying differential
module of the established Nori-to-Hodge realization.
\item It is $t$-exact for the perverse and holonomic hearts and is exact
faithful on those hearts. It commutes with ordinary geometric nearby and
vanishing cycles with their canonical and variation maps.
\item Its absolute restriction is $d_k$, with the actual absolute
comparison, not a newly chosen fiber functor.
\end{enumerate}
The same statements hold after the $k$-linear coefficient extension of the
rational Nori category.
\end{theorem}
\begin{proof}
\Cref{lem:module-extension} gives the functor, its tensor structure, and
its formula on free modules. Since the augmentation is pulled back from
$\Spec k$, associativity and pullback of relative tensor products supply
the pullback comparison. On the absolute category the bar lemma identifies
the formula with $d_k$: by construction \eqref{eq:deRham-augmentation}
is exactly its action of $r_k(\mathscr N_k)$ on the unit.

For the complex comparison, write the established Betti functor on
$\mathsf{DN}(X)$ in the module model. It is linear over the Betti
realization of $\mathsf{DM}(X)$; its action on the unit is its own
augmentation, pulled back from $\Spec k$. The absolute comparison and
\cref{lem:absolute-dR-comparison} identify this augmentation with
$\alpha_k\otimes_k\C$ under the comparison for $r_k$. Pullback identifies
the actions over $X$. The bar lemma now identifies the two functors on
\emph{all} modules and their morphisms. This proves (ii) as a natural
enhanced comparison, rather than by assigning unrelated descent data to
the Hodge realizations.

We give the six-operation verification explicitly. The compact module
category is the thick closure of the free modules on compact geometric
motives. Indeed, these modules are compact by the free--forgetful
adjunction, since the forgetful functor preserves filtered colimits; they
generate, since a module with no maps from their shifts has zero underlying
object. This is also the compact generation used in Tubach's \S4.2.
The free modules realize as constructible regular-holonomic objects by
\cref{lem:geometric-dR}; finite cones and retracts show that the realization
of every compact module is bounded constructible.

The exchange map for a direct image is the mate of the pullback
comparison and the adjunction counit. On a free module
$\mathsf{Nor}_Y^*Q$, the source uses
$f_*\mathsf{Nor}_Y^*Q\simeq\mathsf{Nor}_X^*f_*Q$ from
\cref{input:ordinary}, and the exchange map becomes exactly
$r_Xf_*Q\to f_*r_YQ$. It is the geometric exchange map, not merely an
isomorphism between its source and target: the free action is multiplication
followed by $\alpha_X$, so naturality of the counit and the pulled-back
augmentation identifies the two mates. This map is an equivalence by
\cref{lem:geometric-dR}. The class of constructible objects on which an
exchange map between exact functors is an equivalence is thick, so the
claim follows for every compact module. The same argument applies to
$f_!$ and the mate for $f^!$. Alternatively one obtains $f_!$ from open
extension, proper direct image, and compactification. Duality is treated
with the evaluation and dualizing object; internal Hom follows from
constructible duality and tensor product. Thus all six comparisons exist
and are isomorphisms on the category used here.

Their coherences are those of the module tensor construction and
adjunction mates. In particular, a composition of pullbacks is identified
by the associativity of the relative tensor product, and taking mates
respects compositions. Projection and base-change squares reduce on free
modules to the geometric squares and are compatible with the bar
construction. We do not deduce equality of derived morphisms merely from
conservativity of Betti realization.

After extension to $\C$, (ii) makes every perverse cohomology outside
degree zero vanish for a Nori heart object. Faithful flatness of
$k\hookrightarrow\C$ reflects that vanishing, proving $t$-exactness.
If a heart object has zero algebraic realization, its Betti realization
is zero; faithful Betti realization makes the object zero. An exact
functor detecting zero objects is faithful by applying this observation
to images of morphisms. This also identifies our heart functor with the
canonical universal-abelian-hull factorization of
$Q\mapsto{}^pH^0r_X(Q)$: after comparison its kernel is the same as that
of perverse Betti realization, and the universal factorization is unique.

For ordinary geometric cycles one can use the actual gluing diagrams.
Unipotent cycles are stabilized kernels and cokernels formed from
$j_!,j_*$ and the logarithm objects whose algebraic connections are the
nilpotent logarithmic systems $d-N\,dt/t$. Their units, trace maps, and
nilpotent endomorphisms are geometric over $k$. The functor just constructed
preserves these diagrams and their maps. Finite Kummer pullbacks and their
character projectors supply the other eigenvalues for geometric objects;
over $k=\overline\Q$ those projectors are defined. This proves the ordinary
nearby/vanishing-cycle assertion including the canonical and variation
maps. It does not identify arbitrary \emph{irregular} moderate cycles with
topological cycles; that separate issue is treated in
\cref{prop:regular-cycle-bridge}.

Finally extend coefficients, not the geometric base. For a finite
coefficient field $E/\Q$ in $k$, first apply the rational realization
with its $E$-action, then take the summand corresponding to the specified
embedding $E\hookrightarrow k$ in $E\otimes_\Q k$. Separability makes
that projector exact. Every object, morphism, and finite diagram in the
$k$-linear scalar extension is defined over such a finite $E$ and then
idempotent completion; all comparisons above commute with these
projectors. Passing to their filtered union gives the final assertion.
This is distinct from the geometric extension $k\hookrightarrow\C$.
\end{proof}

\begin{remark}[What has been discharged]\label{rem:interface-discharge}
The relative algebraic functor, its operation maps, and the comparison are
provided by \eqref{eq:deRham-augmentation}--\eqref{eq:relative-algebraic-realization}.
The relative tensor product in the formula is genuinely derived.
Neither its $t$-exactness nor its faithfulness is automatic for a general
augmented algebra; those conclusions were proved by the actual complex
comparison and faithful scalar extension. Likewise, agreement on the
objects of a generating family alone would not suffice without agreement
on the augmentation and the bar-resolution maps.
The standard external inputs are now the ordinary Nori module presentation,
the ordinary geometric differential-module formalism, and the absolute
mixed-realization construction. No relative algebraic realization theorem
or relative period theorem is inserted as an assumption in their place.
The construction is deliberately on enhanced categories: a statement only
about isolated cohomology groups would not define the augmentation and
bar comparison. A new independent audit may challenge these proofs or an
application of a cited input; the description here records the explicit
replacement argument, not an independent certification.
\end{remark}

\subsection{The affine-line quotient}

Fix $X$ and $\pi:X\times\A^1\to X$. Let
$\mathcal A_X=\MN(X)$ and
$\Ccal_X=\pi^*[1]\mathcal A_X\subset\mathcal A_{X\times\A^1}$.

\begin{lemma}[Derived affine-line full faithfulness]\label{lem:affine-ff}
The adjunction unit $\id\to R\pi_*\pi^*$ is an isomorphism. In particular
$\pi^*$ is fully faithful on the derived categories, in all degrees.
\end{lemma}
\begin{proof}
Under Betti realization the unit is the usual pullback--pushforward unit
for the product with the contractible complex affine line, and is an
isomorphism. The cone has zero derived Betti realization; conservativity
makes the cone zero. Adjunction now gives
\[
 \Hom(\pi^*A,\pi^*B[n])=\Hom(A,B[n])
\]
for every $n$. This is an assertion about the actual adjunction unit, not
merely an unspecified equivalence of objects.
\end{proof}

\begin{lemma}[Reflection of pulled-back heart objects]\label{lem:ordinary-reflection}
If $A\in\mathcal A_{X\times\A^1}$ and
$\rat(A)\simeq\pi^*L[1]$ for a perverse sheaf $L$ on $X$, then
$A\simeq\pi^*B[1]$, where $B=\pH^{-1}(R\pi_*A)$.
\end{lemma}
\begin{proof}
The Betti realization of $R\pi_*A$ is $L[1]$. Its other perverse
cohomology groups vanish, and exact faithfulness reflects their vanishing.
Thus $R\pi_*A\simeq B[1]$. The counit
$\pi^*B[1]\to A$ realizes as an isomorphism. Its cone is zero by
conservativity.
\end{proof}

\begin{proposition}[Serre closure and the derived essential image]\label{prop:serre-derived}
The subcategory $\Ccal_X$ is Serre. The essential image of derived
$\pi^*$ consists precisely of the bounded complexes whose perverse
cohomology belongs to $\Ccal_X$. It is thick and stable under perverse
truncations.
\end{proposition}
\begin{proof}
Smooth pullback with connected fibers has subquotient-closed image on
perverse sheaves, by BBD, Corollaire~4.2.6.2. Exact Betti realization and
\cref{lem:ordinary-reflection} give subquotient closure in the Nori heart.

For extension closure, an extension of $\pi^*B[1]$ by $\pi^*A[1]$ is
represented by a connecting morphism
$\pi^*B[1]\to\pi^*A[2]$. By \cref{lem:affine-ff} this is the pullback
of a morphism $B\to A[1]$. Its extension triangle is a pullback triangle,
and $\pi^*[1]$ is $t$-exact; hence the middle heart object belongs to
$\Ccal_X$. This proves the Serre property using the special affine-line
adjunction, not a claim about arbitrary connected smooth fibers.

The pullback of a bounded complex has cohomology in $\Ccal_X$, up to the
fixed shift. Conversely, induct on the number of nonzero perverse
cohomology objects of a complex $K$. A single-cohomology object is in the
image. In a truncation triangle, the two shorter objects are in the image
by induction. Derived full faithfulness lifts the connecting morphism, so
its cone is again a pullback. This proves the essential-image description.
It also gives truncation stability. A summand is in the image because its
cohomology is a summand of the pulled-back cohomology, hence belongs to the
Serre subcategory. Thus the image is thick.
\end{proof}

\begin{remark}[Two invalid general principles are not used]
The pullback image for $\Gm\to\Spec k$ is not extension-closed: the
connection $d-\left(\begin{smallmatrix}0&1\\0&0\end{smallmatrix}\right)
dx/x$ is a nonconstant extension of two trivial connections. Likewise a
Serre inclusion need not induce a fully faithful bounded derived functor.
The proof above uses derived affine-line full faithfulness before
truncation induction. It agrees with the division of responsibilities in
\cite[Proposition~7.2 and \S7.3]{Virk}.
\end{remark}

\begin{theorem}[The quotient heart and faithful realization]\label{thm:quotient}
The Verdier quotient
\[
 D^b(\mathcal A_{X\times\A^1})/\pi^*D^b(\mathcal A_X)
\]
has the induced bounded $t$-structure with heart
$\EM(X)=\mathcal A_{X\times\A^1}/\Ccal_X$.
Exponentiation supplies the derived six operations and derived convolution.
Hodge realization descends to an exact faithful functor on the quotient
heart. Composing with nonzero exponential realization gives an exact
faithful $\rho_\lambda$, which reflects the bounded $t$-structure.
\end{theorem}
\begin{proof}
By \cref{prop:serre-derived}, the localization subcategory is stable under
truncations. The localization theorem for $t$-structures identifies the
heart with the indicated Serre quotient. The derived operations are those
of exponentiation \cite{GallauerPepin}; this statement does not claim
that all six operations are exact on every heart.

The Hodge functor maps the constant Serre subcategory into the corresponding
Hodge subcategory. It thus induces an exact quotient functor. If a quotient
object has zero Hodge image, any heart representative has Hodge realization
pulled back from $X$, hence Betti realization pulled back from $X$.
\Cref{lem:ordinary-reflection} puts the representative in $\Ccal_X$.
The induced exact functor therefore detects zero objects, and is faithful:
apply detection to the image of any morphism it kills. This argument does
not require an arbitrary quotient morphism to lift to a morphism between
particular chosen representatives.

Virk's Theorems~7.5--7.6 give exactness and faithfulness of nonzero
realization on the exponential Hodge heart. Composition gives the same for
$\rho_\lambda$. Finally
$\rho_\lambda(\pH^iK)=H^i(\rho_\lambda K)$, so exact faithfulness detects
all cohomology vanishings and reflects the bounded $t$-structure.
\end{proof}

\subsection{The algebraic exponential functor and the matching Betti fiber}

\begin{theorem}[Algebraic nonzero-scale realization]
\label{thm:algebraic-exponential}
For $\lambda\in k^\times$ and an ordinary representative
$A\in D^b(\MN(S\times\A^1))$, define
\begin{equation}\label{eq:algebraic-exponential-formula}
 \rho_{\lambda,k}(\mathsf Q_S A)
 =p_!\bigl(R_{\dR,S\times\A^1}(A)
             \otimes\mu^*\mathscr L_{\lambda,k}\bigr),
 \qquad \mathscr L_{\lambda,k}=\mathscr K_{\lambda,k}[-1].
\end{equation}
Here $p$ and $\mu$ are the two projections and
$\mathscr K_{\lambda,k}=(\Ocal_{\A^1_k},d-\lambda\,dt)$ is the
holonomic heart object. The formula is a well-defined algebraic functor
on the exponential derived quotient. It is tensor-compatible, commutes
with $f^*,f_*,f_!,f^!$, and has duality
\[
 \DD\rho_{\lambda,k}=\rho_{-\lambda,k}\DD.
\]
It is $t$-exact and faithful on the exponential heart. Its complexification
is Virk's realization of the Hodge realization, with its actual operation
maps. We henceforth denote this algebraic functor by $\rho_\lambda$.
\end{theorem}
\begin{proof}
All terms in \eqref{eq:algebraic-exponential-formula} are over $k$ by
\cref{thm:algebraic-interface}. The polynomial operator
$\partial_t-\lambda$ is bijective on $k[t]$: on a polynomial $v$ of
degree $m$ its inverse is
\[
 -\sum_{j=0}^{m}\lambda^{-j-1}\partial_t^jv.
\]
Thus the algebraic direct image of the exponential kernel to a point
vanishes; by duality the compactly supported direct image vanishes too.
Projection formula now shows that the displayed functor kills every
$p^*$-object. It therefore factors through the Verdier quotient, including
all its morphisms, rather than merely being independent of representatives
at the level of object classes.

The identities
\[
 \mathrm{sum}^*\mathscr L_{\lambda,k}
 \simeq\mathscr L_{\lambda,k}\boxtimes\mathscr L_{\lambda,k},
 \qquad i_0^*\mathscr L_{\lambda,k}=k
\]
are the addition and zero-section identities for the shifted exponential
kernel. Projection formula and Kunneth give tensor compatibility and the
unit, with their associativity constraints; see also Virk's
Propositions~6.4 and~6.8. Proper-support composition and base change first
give $f_!$ and $f^*$ compatibility. The natural forget-support map
$p_!\to p_*$ after the twist is an isomorphism on ordinary regular
input. Indeed, after $k\hookrightarrow\C$ it is exactly Virk's
forget-support isomorphism, proved in his Appendix~B. Algebraic base
extension is faithfully flat and commutes with the transfer complexes, so
it reflects this isomorphism. Duality changes the sign of the exponential
kernel, giving the displayed scale-reversing formula and the remaining
$f_*,f^!$ compatibilities. These are the same projection, support, and
adjunction maps as after complexification.

By \cref{thm:algebraic-interface}, the complexified ordinary input is
the underlying module of its Hodge realization. Thus the complexification
of \eqref{eq:algebraic-exponential-formula} is Virk's actual formula, not
an arbitrary isomorphic connection. Virk's Theorems~7.5--7.6 prove
$t$-exactness and heart faithfulness there. Faithful flatness reflects
cohomology vanishings over $k$, and the faithful Hodge quotient realization
of \cref{thm:quotient} reflects zero objects. Exactness then detects
images of morphisms and proves faithfulness of $\rho_{\lambda,k}$.
\end{proof}

\begin{proposition}[Comparison with the chosen rapid-decay fiber]
\label{prop:relative-Betti-interface}
Choose the direction used in \eqref{eq:betti-fiber} consistently with
$e^{-\lambda t}$, and let $s_0\in S(k)$ be a lisse point. The algebraic
fiber $\omega_{v}(N)=\rho_\lambda(N)_{s_0}$ and the chosen Betti fiber
have a natural tensor comparison
\[
 c_N:\omega_{\B,s_0}(N)\otimes_k\C
       \xrightarrow{\sim}\omega_v(N)\otimes_k\C.
\]
On a sufficiently small simply connected lisse analytic open $B$ about
$s_0$, these comparisons extend uniquely to a horizontal tensor
comparison with the analytified algebraic connection. On elementary
exponential data this is the usual rapid-decay integral comparison.
\end{proposition}
\begin{proof}
Normalized restriction $i_{s_0}^*[-d]$ commutes with the realization just
constructed. At a point its de Rham functor is the absolute exponential
de Rham functor: the ordinary part is $d_k$ by
\cref{thm:algebraic-interface}(iv), and the exponential part is its
twisted algebraic direct-image formula. The directional nearby-fiber
functor at infinity and the absolute de Rham--rapid-decay comparison
therefore give $c_N$, functorially in the absolute exponential object.
Equivalently, compute on a regular representative on the potential line:
the de Rham--rapid-decay theorem and the directional nearby-fiber
identification give the comparison of its two vector spaces. The maps
are natural in that representative and kill the constant subcategory, so
they descend to the quotient. This uses the period comparison of these
realizations, not an unproved equivalence of two entire absolute motivic
categories.
These are the absolute comparison inputs of
\cite{FresanJossen,Hien}; the directional convention and tensor fiber
are recorded in \cite[Definition~3.5 and Theorem~3.6]{Snodgrass}.
Cohomological variance is used on both sides; a homology pairing is
converted to this comparison by its duality, not by identifying a homology
space with itself as a cohomology space.

Parallel transport of a flat connection extends each $c_N$ uniquely over
$B$. For a morphism $N\to M$, the two resulting maps of connections are
horizontal and agree at $s_0$, hence agree on $B$. The same argument for
the tensor, symmetry, unit, and evaluation maps proves the tensor
coherences. It applies to all subquotients because the absolute comparisons
are already natural on the absolute exponential category and the
algebraic realization is exact. Thus we are not inferring naturality from
objectwise choices of a period matrix.

For elementary data, rapid-decay integrals are horizontal solutions of
the twisted Gauss--Manin connection and their values at $s_0$ are the
absolute comparison just used \cite{HienRoucairol,Hien}. Uniqueness of
horizontal continuation makes the two comparisons equal on $B$. In
particular the later identity $P=Yc$ uses this specific comparison, with
no algebraicity assertion for the entries of $c$.
\end{proof}

\section{Fixed parts and extension from dense opens}\label{sec:fixed}

All tensor operations in this section are normalized as in
\eqref{eq:lisse-tensor}. Call $N\in\EM(S)$ lisse when $\rho_\lambda(N)$
is an algebraic vector bundle with integrable connection. Exactness,
faithfulness, and the usual closure properties of algebraic connections
make this a Serre subcategory. Work in its dualizable rigid envelope when
using Tannakian language.

\begin{lemma}[Lowest direct image and independent horizontal sections]\label{lem:lowest}
For a connection $M$ on a smooth connected $d$-fold,
\[
 H^j(a_+M)=0\ (j<-d),\qquad
 H^{-d}(a_+M)=H^0_{\dR}(S,M)=\Gamma(S,M)^\nabla.
\]
Evaluation $\Ocal_S\otimes_k\Gamma(S,M)^\nabla\to M$ is injective and
its image is the maximal trivial algebraic subconnection.
\end{lemma}
\begin{proof}
Algebraic de Rham hypercohomology has no negative ordinary degree; the
perverse shift is $[d]$. Its lowest group is the kernel of the degree-zero
connection operator on global sections. This gives the displayed formula
also for a nonaffine base.

Horizontal sections independent over $k$ are independent over $K=k(S)$.
Indeed, in a shortest nontrivial rational relation normalize one coefficient
to one and differentiate. Minimality forces every remaining coefficient to
be a common differential constant, hence to lie in $k$. This contradicts
$k$-independence. Locally analytically, a horizontal section vanishing at
one point vanishes identically; equivalently a fundamental solution
matrix makes independent horizontal sections a constant full-rank matrix.
Evaluation is therefore a bundle injection. Every trivial subconnection
has a basis of global algebraic horizontal sections, proving maximality.
\end{proof}

\begin{theorem}[Algebraic exponential fixed part]\label{thm:fixed}
Let $N$ be lisse and set
\[
 C_N=\pH^{-d}(a_*N),\qquad
 N^{\fix}=\im\bigl(a^*C_N[d]\longrightarrow N\bigr).
\]
Then $\rho_\lambda(N^{\fix})$ is the maximal trivial algebraic
subconnection of $\rho_\lambda(N)$. An object has trivial connection
realization if and only if it is constant, namely of the form $a^*C[d]$.
The object $N^{\fix}$ is the maximal constant subobject of $N$.
\end{theorem}
\begin{proof}
Operation compatibility and \cref{lem:lowest} make the perverse cohomology
of $a_*N$ vanish below $-d$ after realization. Reflection gives that
vanishing in the motivic category. The lowest truncation map and adjunction
therefore define the displayed morphism. It realizes as horizontal
evaluation. Exact realization takes its image to the maximal trivial
subconnection.

If $Q$ realizes a trivial connection, its own evaluation map
$a^*C_Q[d]\to Q$ realizes an isomorphism, and hence is an isomorphism.
Conversely a constant object has constant connection realization. Apply
this to $N^{\fix}$. A morphism $a^*C[d]\to N$ corresponds by adjunction
to $C[d]\to a_*N$. The lowest-cohomology truncation and the $t$-structure
factor this map through $C_N[d]$. Thus every constant subobject is
contained in $N^{\fix}$.
\end{proof}

\begin{corollary}[Constants and their fiber functors]\label{cor:constants}
Constant lisse objects are closed under subquotients, finite direct sums,
normalized tensor products, and tensor duals. The functor $a^*[d]$ is
fully faithful on the absolute heart. On the constant subcategory, the
generic de Rham fiber is scalar extension from a $k$-valued absolute
fiber functor, and its actual comparison matrix is independent of $s$.
\end{corollary}
\begin{proof}
Subquotients of a trivial connection are trivial, so use
\cref{thm:fixed}. The other closure properties follow in the same way or
from normalized pullback. For full faithfulness, adjunction computes
$\Hom(a^*C[d],a^*D[d])$ as $\Hom(C,a_*a^*D)$. By the projection formula,
$a_*a^*D=D\otimes a_*\one_S$. Connectedness and the degree-zero direct
image give $H^0(a_*\one_S)=\one$ and no negative cohomology. Maps from
$C$ see just that degree-zero term, by truncation orthogonality. This gives
$\Hom(C,D)$. Pullback and comparison for the absolute object now give
the fiber-functor statement.
\end{proof}

\begin{remark}[Algebraic, not merely analytic, invariants]
On $\A^1_s$, $(\Ocal,d-ds)$ has the analytic horizontal section $e^s$
but no nonzero algebraic horizontal section: $p'=p$ has no nonzero
polynomial solution. The fixed-part theorem takes algebraic horizontals.
Finite Betti monodromy alone does not make an algebraic connection constant.
\end{remark}

Let $j:U\hookrightarrow S$ be a dense open immersion. Intermediate
extension in the exponential heart is the image of
$\pH^0j_!\to\pH^0j_*$; its formation commutes with exact realization.

\begin{lemma}[Lisse intermediate extension]\label{lem:middle}
A lisse exponential object $N$ is canonically $j_{!*}j^*N$. Restriction
between the lisse subcategories is fully faithful. If $Q_U\in\EM(U)$
realizes $L|_U$ for a connection $L$ lisse on $S$, then
$Q=j_{!*}Q_U$ is lisse and realizes $L$.
\end{lemma}
\begin{proof}
A lisse connection on $S$ has no nonzero holonomic subobject or quotient
supported in $S\setminus U$, and is the minimal extension of its
restriction. Exact realization reflects the absence of such boundary
objects. The natural adjunction morphisms therefore identify $N$ with
$j_{!*}j^*N$.

A morphism on $U$ induces a morphism of intermediate extensions. The
canonical identifications just proved extend it to $S$; a morphism with
zero restriction has boundary-supported image, hence is zero. This proves
full faithfulness. Finally
$\rho_\lambda(j_{!*}Q_U)=j_{!*}(L|_U)=L$. Reflection places $Q$ in the
lisse heart.
\end{proof}

\begin{proposition}[Extending invertible lifts]\label{prop:extend-line}
In the last assertion of \cref{lem:middle}, suppose $Q_U$ is invertible.
Then $Q$ is invertible on $S$.
\end{proposition}
\begin{proof}
Extend $Q_U^{-1}$ by intermediate extension to $Q'$. Its realization is
$L^{-1}$, so $Q'$ is lisse. The products $Q\otimes_{\lis}Q'$ and
$Q'\otimes_{\lis}Q$ are lisse and restrict to the unit on $U$. Applying
\cref{lem:middle} to these products and to the unit identifies both
products with $\one_{S,\lis}$. These isomorphisms extend the inverse
isomorphisms on $U$ and satisfy their coherence identities by faithful
restriction.
\end{proof}

\begin{remark}[What this does and does not say about opens]
Every generic differential subquotient of a fixed lisse connection has a
unique lisse model on $S$: take the saturated differential subobject in
the holonomic category, then its quotient. Coherent modules with integrable
connection on a smooth characteristic-zero variety are locally free, so
no interior singularity is introduced. Morphisms extend by the same
minimal-extension property. The preceding proposition allows an
individually constructed motivic rank-one lift to be returned to that
fixed base. It does not assert that the intersection of infinitely many
dense opens is open, or that finite tensor generation means finitely many
tensor constructions. The corresponding lisse-extension and full-faithfulness
arguments apply to finite diagrams in the generic motivic category.
\end{remark}

\section{Fourier fibers and ordinary geometric constituents}\label{sec:fourier}

For this section an ordinary geometric object means one in the ordinary
Nori/Hodge realization framework, before exponentiation. On smooth dense
opens we use the abelian subquotient-closed category of ordinary geometric
local systems; Jacobsen--Terenzi, Corollary~5.15, identifies the two standard
Nori and Ayoub descriptions \cite{JacobsenTerenzi}. No irregular
Riemann--Hilbert correspondence is being used to algebraize a connection.

\begin{externalinput}[The Fourier comparison used here]\label{input:fourier}
For regular geometric input on $S\times\A^1_t$, use the normalized
partial Fourier--Laplace transform. The finite-graph microlocalization
result gives a connection on $S\times\Gm$ after shrinking $S$.
For a pure polarizable minimal extension $\mathcal T$ on
$S\times\mathbb P^1_t$, Sabbah's Proposition~5.8, Corollary~5.20, and
sequences (5.21)--(5.22) compare the Fourier cycles at zero with ordinary
cycles at $t=\infty$ \cite{SabbahMonodromy}. Mixed input is reduced using
its strict weight filtration and exact transform/cycle functors.
\end{externalinput}

\begin{lemma}[Finite graph reduction and nonzero-scale lissity]\label{lem:graphs}
For regular geometric perverse input $K_*$ on $S\times\A^1_t$, after
shrinking $S$ there is a finite etale cover on which all dominating
singular multisections split into pairwise disjoint graphs $t=u_i(s)$.
The localized partial transform is a connection on $S\times\Gm$.
\end{lemma}
\begin{proof}
Take a finite Whitney stratification adapted to the regular input. Delete
the closures of the images of strata that do not dominate $S$. Each
remaining proper stratum has dimension $d$: it dominates a $d$-dimensional
base and lies properly in a $(d+1)$-fold. Its closure is generically finite
over $S$. Normalize these finitely many multisections and shrink so that
they are finite etale. A common finite Galois refinement splits them into
sections. Delete the images of intersections of distinct sections. The
characteristic variety is then contained in the zero section and the
conormals of the resulting graphs.

Apply the parameterized localized Laplace/microlocalization input with
$z=\lambda^{-1}$. In Sabbah's Appendix~B the localized module is locally
free over the parameter ring $\Ocal[z,z^{-1}]$, with integrable
connection; Proposition~B.1 identifies its residual systems. The reference
is \cite[Appendix~B, preprint p.~12, immediately before Proposition~B.1]{SabbahParameters},
not p.~40 of that preprint. The underpinning general parameter statement is
\cite[Propositions~1.18 and~1.20]{DouaiSabbah}. This gives nonzero-scale
lissity locally on the cover, which descends in the finite etale topology.
\end{proof}

\begin{lemma}[Fixed fiber and topological nearby local system]
\label{lem:fiber-nearby}
After the shrinking of \cref{lem:graphs}, a fixed nonzero Fourier fiber
and the \emph{topological} nearby-cycle local system at $\lambda=0$
have isomorphic underlying local systems on $S$, after forgetting the
commuting nearby monodromy.
\end{lemma}
\begin{proof}
The local system on $S^{\mathrm{an}}\times\C^\times$ is a representation
of $\pi_1(S)\times\mathbb Z$. Restriction to a fiber forgets the second
operator. Topological nearby cycles have the same representation of
$\pi_1(S)$ and retain the second operator as nearby monodromy. Choose
a path from the fixed scale to a punctured neighborhood of zero. Changing
its winding changes the isomorphism by a power of the commuting operator,
not the underlying isomorphism class or its constituents.
\end{proof}

\begin{proposition}[The regular-specialization bridge]
\label{prop:regular-cycle-bridge}
For the partial Fourier transform of the regular geometric input used in
\cref{input:fourier}, moderate algebraic nearby cycles at $\lambda=0$
realize the topological nearby cycles in \cref{lem:fiber-nearby}.
Consequently those moderate cycles and a fixed nonzero Fourier fiber have
the same underlying $S$-local-system constituents after forgetting the
nearby-monodromy operator. This assertion uses regular specialization
along $\lambda=0$, not just lissity on $S\times\Gm$.
\end{proposition}
\begin{proof}
For pure input take the regular minimal extension $\mathcal T$ in the
$t$-direction and its Hodge-to-twistor realization. In Sabbah's notation,
the \emph{kernel} before projective direct image is
$\mathcal F\mathcal T=p^+\mathcal T[*\infty]\otimes e^{-t\lambda/z}$.
Its underlying strict module is strictly specializable and regular along
$\lambda=\lambda_0$ for every finite $\lambda_0$, including zero, by
\cite[Proposition~4.1(ii)]{SabbahMonodromy}. Notice the contrast with
Proposition~4.1(i), which allows irregularity along $t'=0$ at nonzero
scale; regularity in the two different directions must not be interchanged.

Use the moderate-cycle/de Rham comparison for a module regular along
the specialization hypersurface, with the perverse shift $[-1]$ on
topological nearby cycles. The regular-specialization hypothesis is what
removes a possible purely irregular summand invisible to moderate cycles.
Proper direct image is compatible with these regular-specialization
cycles and with their de Rham comparison. Apply it to the projective
$t$-projection: the resulting comparison is that of the moderate nearby
cycles of the partial transform with its topological nearby cycles.
These are the regular-specialization and proper-image comparison maps,
not an invocation of regular Riemann--Hilbert for an arbitrary irregular
connection. The cycle comparisons in Sabbah's Proposition~5.8 and
Corollary~5.20 are used in precisely this setting.

For mixed geometric input, use the finite strict weight filtration of
its Hodge realization. Normalized Fourier transformation and both cycle
functors are exact; the cycle comparison is natural. The pure comparison
on each graded constituent therefore gives the comparison for the mixed
object by the finite filtration and the long exact cohomology sequences.
Finally use \cref{lem:fiber-nearby}. All these statements concern the
complex realization; \cref{thm:algebraic-interface} supplies the
algebraic ordinary objects and their geometric operation maps.
\end{proof}

\begin{remark}[The lissity-only assertion would be false]
The connection $\Ecal^{1/\lambda}$ on a punctured disc has a trivial
rank-one analytic horizontal local system, but its formal type is purely
irregular at zero and its moderate nearby cycles vanish. Its topological
nearby local system has rank one. It is therefore essential to insert
\cref{prop:regular-cycle-bridge} before using the topological product
argument in the algebraic-cycle proof. This example does not satisfy the
regular-specialization conclusion for the intended Fourier kernel.
\end{remark}

\begin{proposition}[Unipotent comparison through two triangles]\label{prop:two-triangles}
Let $t'=t^{-1}$ and let $q:S\times\mathbb P^1_t\to S$ denote the
proper projection after taking cycles in the Fourier coordinate. Set
\[
 V=\psi_{\lambda,-1}(\mathcal F\mathcal T),\qquad W=\ker N_\lambda.
\]
Every simple constituent of every perverse cohomology of $q_+V$ occurs
among the perverse-cohomology constituents of the proper direct images of
\[
 i_{\infty,+}\ker N_{t'},\qquad
 \mathcal T,\qquad
 i_{\infty,+}\psi_{t',-1}\mathcal T.
\]
All cycles here are perverse-normalized.
\end{proposition}
\begin{proof}
Before $q_+$, Sabbah's (5.22) and the first sequence of (5.21) give
\begin{align}
0&\longrightarrow i_{\infty,+}\ker N_{t'}
 \longrightarrow W\longrightarrow\mathcal T\longrightarrow0,
 \label{eq:upstream-one}\\
0&\longrightarrow W\longrightarrow V
 \longrightarrow i_{\infty,+}\psi_{t',-1}\mathcal T\longrightarrow0.
 \label{eq:upstream-two}
\end{align}
Those sequences are stated before the subsequent dimension-one
specialization in Sabbah's argument. Apply derived proper direct image to
obtain two distinguished triangles. For any triangle $A\to B\to C\to$,
the segment
$\pH^jA\to\pH^jB\to\pH^jC$ gives a short exact sequence whose middle
object is $\pH^jB$, whose subobject is a quotient of $\pH^jA$, and whose
quotient is a subobject of $\pH^jC$. Apply this observation
first to \eqref{eq:upstream-one}, then to \eqref{eq:upstream-two}, for
each $j$. Finite length gives exactly the asserted constituent containment.
No exactness of $q_+$ on the heart, and no interchange of a constituent
filtration with a derived functor, is asserted.
\end{proof}

\begin{theorem}[Finite Betti monodromy of rank-one factors]\label{thm:finite-betti}
A rank-one differential subquotient, defined over $k$, of a fixed nonzero
Fourier fiber of regular geometric input has finite Betti monodromy on a
suitable dense smooth open.
\end{theorem}
\begin{proof}
Use \cref{prop:regular-cycle-bridge} to identify the moderate-cycle
objects with the topological nearby local systems. For pure input, the
non-unipotent Fourier comparisons identify the relevant
constituents with ordinary nearby-cycle constituents at infinity. The
unipotent case follows from \cref{prop:two-triangles} and proper
compatibility of nearby cycles. Minimal extension, ordinary cycles,
proper direct image, and perverse cohomology remain ordinary geometric
operations. For mixed input, use the strict finite weight filtration on
the Hodge realization, exactness of normalized Fourier transformation and
nearby cycles, and semisimplicity of each pure graded Hodge object. Its
simple constituent local systems belong to the ordinary geometric-origin
category; the abelian subquotient closure and its compatibility with the
Nori description are the comparison input of \cite[Corollary~5.15 and Proposition~5.17]{JacobsenTerenzi}.
Apply the pure argument to the graded constituents.

By \cref{lem:fiber-nearby,prop:regular-cycle-bridge}, a rank-one Betti constituent of the fixed
fiber is one of these ordinary geometric local systems. The local system
of a rank-one differential subquotient is such a constituent, because
analytic horizontal realization is exact on connections. Lam--Litt's
published Lemma~3.5 says that an ordinary geometric rank-one local system
is torsion \cite{LamLitt}. Its image is therefore finite. This gives
finite Betti monodromy, not regular singularity or finite tensor order of
the original algebraic connection.
\end{proof}

\section{The arithmetic kernel and rank-one lifts}\label{sec:arithmetic}

\begin{externalinput}[Degree-one arithmetic realization]\label{input:one-motives}
Use the cohomological Picard $1$-motive
\[
 \operatorname{Pic}^{+}(U)=
 [\operatorname{Div}^{0}_{D}(\overline U)\to
   \operatorname{Pic}^{0}(\overline U)]
\]
for a smooth compactification with normal-crossings boundary. Its realization
is $H^1(U)(1)$. A rational boundary divisor giving a morphism from the
weight-zero lattice object realizes, when made principal, as its logarithmic
de Rham class \cite{BarbieriVialeSrinivas}.
The de Rham--Betti realization of $1$-motives over $\overline\Q$ is fully
faithful with image closed under subquotients
\cite[Theorem~9.14]{HuberWustholz}. The comparison with degree-one Nori
motives is \cite{AyoubBarbieriViale}.
\end{externalinput}

\begin{theorem}[Arithmetic exponential kernel]\label{thm:kernel}
Let $U/k$ be smooth and connected and $L$ a rank-one algebraic integrable
connection over $k$ with finite Betti monodromy. On a dense affine open
$U^\circ$ there are $f\in\Ocal(U^\circ)$ and a finite-order algebraic
rank-one connection $A$ such that
\[
 L|_{U^\circ}\simeq\Ecal^f\otimes A.
\]
\end{theorem}
\begin{proof}
The finite Betti character defines a finite etale cover and its character
idempotent over $k$; algebraically closed extension does not change finite
etale covers in characteristic zero. Let $A_0$ be its finite Artin
connection and replace $L$ by $L_0=L\otimes A_0^{-1}$. On a dense affine
open trivialize its underlying line bundle and write $\nabla=d+\omega$.
Integrability gives $d\omega=0$, and trivial monodromy gives
\[
 \int_\gamma\omega\in2\pi i\mathbb Z
 \quad(\gamma\in H_1(U^{\mathrm{an}},\mathbb Z)).
\]
Define the integral Betti class $\eta$ by division of these periods by
$2\pi i$. Comparison sends $[\omega]$ to $2\pi i\eta$. If $\eta=0$,
comparison implies $[\omega]=0$, and affine algebraic de Rham theory
supplies $F\in\Ocal(U)$ with $\omega=dF$. Then $L_0=\Ecal^{-F}$.

Otherwise the \emph{specified} pair $([\omega],\eta)$ gives a morphism
$\Q(-1)\to H^1_{\dRB}(U)$, with our convention that this Tate period is
$2\pi i$. Twist to obtain a specified morphism
$\Q(0)\to H^1_{\dRB}(U)(1)$. By \cref{input:one-motives}, it lifts to
the specified morphism $\Q(0)\to\operatorname{Pic}^{+}(U)_\Q$.
This is a rational boundary divisor $\Delta$ whose Picard class is zero
in $\operatorname{Pic}^0(\overline U)\otimes\Q$.

Choose $N>0$ clearing its denominators and killing its torsion Picard class.
Then $N\Delta=\operatorname{div}(u)$ for a function with no zeros or
poles on $U$, so $u\in\Ocal(U)^\times$. Compatibility with the specified
morphism gives
\[
 N[\omega]=[d\log u]\quad\hbox{in }H^1_{\dR}(U/k).
\]
Since $U$ is affine, closed one-forms modulo exact global forms compute
this cohomology. Thus $N\omega-d\log u=dF$ for some regular $F$, and
\[
 L_0\simeq\Ecal^{-F/N}\otimes
 (\Ocal_U,d+N^{-1}d\log u).
\]
The $N$th power of the second factor has horizontal algebraic generator
$u^{-1}$ and is trivial. Tensor back by $A_0$. This proves the result
without a relative functional-period theorem.
\end{proof}

\begin{proposition}[Realizing the algebraic factors]\label{prop:realize-line}
At scale $\lambda\in k^\times$, $\Ecal^f\otimes A$ in
\cref{thm:kernel} has an invertible exponential motivic lift on $U^\circ$.
If the connection extends lisse to a given smooth base, so does an
invertible lift.
\end{proposition}
\begin{proof}
Use the graph potential $f/\lambda$ in \cref{lem:graph}. Its inverse
is the graph with the negative potential. A finite-order connection is a
rank-one summand of a finite etale pushforward; over $k$ its character
idempotent defines a rank-one Artin motive. Insert it on the zero graph
and take its normalized tensor product with the exponential graph object.
Both factors are invertible and realize the desired factors. The extension
assertion is \cref{prop:extend-line}.
\end{proof}

\begin{theorem}[Rank-one saturation on a fixed lisse model]\label{thm:saturation}
Let $\Tcal$ be a finitely generated rigid subquotient-closed lisse
subcategory of the generic exponential Nori envelope, modeled on $S$.
Suppose its objects have regular geometric pre-Fourier representatives,
as supplied by the exponential quotient and the ordinary geometric
coefficient system. Let $\Dcal$ be the differential tensor category
generated over $K$ by their realizations and differential subquotients.
Every invertible object of $\Dcal$ has an ambient invertible exponential
motivic lift. The lift can be modeled on $S$ whenever the differential
object is modeled lisse there.
\end{theorem}
\begin{proof}
An invertible object $L\in\Dcal$ is a differential subquotient of a
finite tensor construction in realized generators. Choose its motivic
counterpart $B$ and a common open model. The quotient description is
essentially surjective on the heart, so choose an ordinary geometric
representative $K_B$ for $B$. At nonzero scale the representative realizes
as its partial Fourier fiber. Ordinary tensor constructions, convolution,
and the six operations stay within regular geometric input before
exponentiation; no claim of regularity of the exponential output is made.

Apply \cref{thm:finite-betti} to $L$ after shrinking. It has finite
Betti monodromy there. It is algebraic over $k$, so \cref{thm:kernel}
applies. \Cref{prop:realize-line} gives the ambient invertible lift on a
dense affine open, and \cref{prop:extend-line} returns it to $S$ if needed.
This is an object-by-object argument on one fixed lisse model, not an
intersection of infinitely many opens. Rank-one objects with transcendental
connection coefficients are not fed into the arithmetic theorem; extension
of differential constants is handled later in \cref{sec:pv-base-change}.
\end{proof}

\begin{remark}[The optional split Fourier-geometric route]
If a separate construction supplies an actual isomorphism
$q^*L\simeq\Ecal^g\otimes A$ with finite $A$ on a finite cover and an
effective motivic descent datum, its lift is obtained from the potential
$g/\lambda$ and the Artin idempotent, followed by descent and
\cref{prop:extend-line}. Such a statement requires the actual global
rank-one comparison and effective descent hypotheses. A formal Stokes
expression alone is not substituted for them. This optional route is not
used in \cref{thm:saturation}; the full nonzero-fiber argument above is
what supplies saturation for the generic envelope.
\end{remark}

\section{Compatible fibers and Galois exactness}\label{sec:galois}

Choose $s_0\in S(k)$ where a tensor generator is regular and fix the
algebraic de Rham evaluation fiber
\[
 \omega_v:\Dcal\longrightarrow\VecCat_k,\qquad M\longmapsto M_{s_0}.
\]
All objects of $\Dcal$ have their unique lisse model on $S$, as explained
after \cref{prop:extend-line}. The same functor composed with realization
is a fiber functor on $\Tcal$. Set
\[
 G_v=\Aut^\otimes(\omega_v\rho_\lambda|_{\Tcal}),\qquad
 H_v=\Aut^\otimes(\omega_v|_{\Dcal}).
\]
The differential category is generated under subquotients by the image of
$\Tcal$, so the induced morphism $H_v\to G_v$ is a closed immersion,
by \cite[Proposition~2.21(b)]{DeligneMilne}. This is not a consequence
of exact faithfulness alone; the generation statement is also used.

\begin{proposition}[Motivic semi-invariant subspaces]\label{prop:semi}
For $B\in\Tcal$ and a character $\chi:H_v\to\Gm$, the
$\chi$-semi-invariant subspace of $\omega_v\rho_\lambda(B)$ is the
fiber of a subobject of $B$ in $\Tcal$.
\end{proposition}
\begin{proof}
The character defines an invertible differential object $L_\chi$.
By \cref{thm:saturation}, choose an ambient invertible lift
$\mathcal L_\chi$, on a common lisse base if necessary. Take
\[
 C_\chi=(B\otimes_{\lis}\mathcal L_\chi^{-1})^{\fix}.
\]
By \cref{thm:fixed} its connection realization is the maximal trivial
subconnection of $\rho_\lambda(B)\otimes L_\chi^{-1}$. Tensor its
inclusion by $\mathcal L_\chi$ and take the image $W_\chi\subset B$.
Its fiber is exactly
\[
 L_\chi\otimes
 (\rho_\lambda(B)\otimes L_\chi^{-1})^{H_v},
\]
namely the desired semi-invariant subspace. Although the lift was ambient,
$W_\chi$ belongs to $\Tcal$ because that category is closed under
subobjects of its objects.
\end{proof}

\begin{lemma}[Normality from semi-invariants]\label{lem:normality}
Let $H\subset G$ be algebraic groups over an algebraically closed field of
characteristic zero. If $V^\chi$ is $G$-stable for every $G$-representation
$V$ and every character $\chi$ of $H$, then $H$ is normal in $G$.
\end{lemma}
\begin{proof}
By the stabilizer-of-a-line theorem for linear algebraic groups, choose a
$G$-representation $V$ and a line $\ell$ with stabilizer $H$. The line
has an $H$-character $\chi$ and lies in $W=V^\chi$. By hypothesis $W$
is $G$-stable and $H$ acts on it by scalars. For $g\in G$ and $h\in H$,
the element $g^{-1}hg$ therefore preserves $\ell$, so it lies in $H$.
Apply the same argument to $g^{-1}$ to get equality. This proves normality
on geometric points, hence on the characteristic-zero algebraic groups.
This is the semi-invariant form of Andr\'e's criterion used in
\cite{Andre,Mostaed}.
\end{proof}

\begin{theorem}[Exactness in algebraic evaluation coordinates]\label{thm:exact}
Let $\Tcal_0$ be the full constant subcategory of $\Tcal$, and let
$G_{0,v}$ be its group under the evaluation fiber. Then
\[
 1\longrightarrow H_v\longrightarrow G_v
 \longrightarrow G_{0,v}\longrightarrow1
\]
is exact.
\end{theorem}
\begin{proof}
\Cref{prop:semi,lem:normality} make $H_v$ normal. By
\cref{cor:constants}, $\Tcal_0$ is a full tensor subcategory closed
under subquotients. Its inclusion gives a faithfully flat quotient of
$G_v$ \cite[Proposition~2.21(a)]{DeligneMilne}. A $\Tcal$-object is
trivial as an $H_v$-representation precisely when its connection is
trivial. By \cref{thm:fixed} this is equivalent to being constant.
The constant subcategory is also closed under tensor duals, since the dual
of a trivial connection is trivial. Thus the representations of $G_v/H_v$ are precisely $\Tcal_0$, proving
exactness. No period-map injectivity has been used.
\end{proof}

\subsection{Changing to Betti coordinates}

Let $\omega_\B=\omega_{\B,s_0}$ and let $\omega_\dR$ be the generic
$K$-valued de Rham fiber. Define over $k$
\begin{equation}\label{eq:betti-groups}
 G_\B=\Aut^\otimes(\omega_\B|_\Tcal),\qquad
 G_{0,\B}=\Aut^\otimes(\omega_\B|_{\Tcal_0}),\qquad
 H_\B=\ker(G_\B\to G_{0,\B}).
\end{equation}
Let $c:\omega_\B\otimes_k\C\xrightarrow{\sim}
\omega_v\rho_\lambda\otimes_k\C$ be the actual comparison at $s_0$.
It is an isomorphism of fiber functors, not an assertion that its entries
are algebraic.

\begin{lemma}[The group change is conjugation]\label{lem:conjugation}
The comparison $c$ identifies $(G_\B)_\C$ with $(G_v)_\C$, and
\[
 (H_\B)_\C=c^{-1}(H_v)_\C c.
\]
\end{lemma}
\begin{proof}
An isomorphism of fiber functors conjugates their tensor automorphism
functors. Its restriction to $\Tcal_0$ conjugates the quotient groups
as well. Kernels are preserved by this commutative diagram. Use
\cref{thm:exact} to identify the evaluation-coordinate kernel with
$H_v$. This proves the formula before any analytic torsor comparison.
\end{proof}

\section{Normalized solutions and extension of constants}\label{sec:pv-base-change}

The next argument establishes the differential base-change statement
needed in the last section. It prevents arbitrary complex differential
subobjects from being passed unexamined to an arithmetic theorem over $k$.

Shrink around $s_0$ to choose algebraic etale coordinates
$x_1,\ldots,x_d$ and a frame of a tensor-generator connection. The
corresponding derivations $\partial_i$ extend uniquely from
$k(x_1,\ldots,x_d)$ to $K$ and commute. They span the tangent derivations
over $K$. Write the system as
\begin{equation}\label{eq:system}
 \partial_iY=A_iY,\qquad A_i\in M_n(K),\qquad Y(s_0)=I.
\end{equation}
Here $\nabla=d-\sum_iA_i\,dx_i$; integrability is
$\partial_iA_j-\partial_jA_i-[A_i,A_j]=0$.

\begin{lemma}[Algebraic Taylor normalization]\label{lem:taylor}
The solution $Y$ of \eqref{eq:system} exists on a sufficiently small
analytic polydisc $B$ and has Taylor coefficients in $k$.
\end{lemma}
\begin{proof}
In etale coordinates, the completed local ring at $s_0$ is
$k[[x_1,\ldots,x_d]]$, with the coordinates translated to zero.
The $A_i$ have coefficients there. Coefficient comparison in
\eqref{eq:system}, starting from $I$, determines each Taylor coefficient
using only coefficients already determined and division by positive
integers. Integrability is precisely the consistency of the different
recursions. Analytic existence for a nonsingular flat system gives a
convergent solution; uniqueness identifies its Taylor series with this
formal solution. Thus all its coefficients lie in $k$.
\end{proof}

\begin{externalinput}[Picard--Vessiot facts for integrable systems]\label{input:pv}
For finitely many commuting derivations and an algebraically closed common
constant field, a field generated by a fundamental matrix with no new
common constants is Picard--Vessiot; its fundamental-matrix algebra is
differentially simple. Its spectrum is the differential comparison torsor,
and its group is the tensor automorphism group of the horizontal fiber
functor. Use classical Picard--Vessiot theory and its integrable-system
form \cite{PutSinger,Kolchin,CHM}. This is ordinary Picard--Vessiot theory
with all base derivations, not parameterized Picard--Vessiot theory with
some derivatives omitted.
\end{externalinput}

Put
\[
 R_k=K[Y,\det(Y)^{-1}]\subset\mathscr M(B),\qquad E=\C(S).
\]
The fraction field of $R_k$ has common constants $k$. Indeed, a horizontal
meromorphic germ is complex constant; after writing the germ as a quotient
of two formal series over $k$, a nonzero coefficient of the denominator
shows that this constant belongs to $k$. Thus \cref{input:pv} applies to
$R_k$, and its solution fiber is the evaluation fiber $\omega_v$, since
$Y(s_0)=I$. Its group is $H_v$.

\begin{lemma}[Finite coefficient descent]\label{lem:finite-linear}
For $f_1,\ldots,f_r\in k[[x_1,\ldots,x_d]]$, the space of complex
constant relations among the $f_j$ is the complex scalar extension of
the space of $k$-relations.
\end{lemma}
\begin{proof}
Expand $f_j=\sum_\alpha a_{\alpha j}x^\alpha$. A relation is the
kernel of the matrix $(a_{\alpha j})$ with $r$ columns. Its possibly
infinite family of rows spans a finite-dimensional subspace of $k^r$.
Choose finitely many rows that are a basis of this row span. They define
the same kernel, and kernels of finite matrices commute with field
extension. This proves the claim without a field automorphism or a
continuity assumption on one.
\end{proof}

\begin{proposition}[No extra relations after complex constant extension]\label{prop:pv-base-change}
Multiplication of meromorphic functions gives an isomorphism
\begin{equation}\label{eq:pv-extension}
 R_k\otimes_K E\xrightarrow{\sim}
 E[Y,\det(Y)^{-1}]=:R_E.
\end{equation}
The algebra $R_E$ is Picard--Vessiot over $E$ with group $(H_v)_\C$.
\end{proposition}
\begin{proof}
Surjectivity is immediate. To prove injectivity, take a finite expression
in the kernel. Choose an affine neighborhood of $s_0$ defined over $k$.
Every coefficient in $E$ is a quotient of finite complex linear
combinations of regular $k$-functions. Clear a common nonzero denominator
of this type. The alleged relation becomes a finite expression
\[
 \sum_{j=1}^r c_j r_j=0,\qquad c_j\in\C,\quad r_j\in R_k.
\]
Multiply once more by a nonzero element of $K$, if necessary, to remove
the finitely many poles at $s_0$ of the $r_j$. Their expansions are now in
$k[[x_1,\ldots,x_d]]$ by \cref{lem:taylor}. By
\cref{lem:finite-linear}, the vector $(c_j)$ is a complex linear
combination of $k$-relation vectors. Those relations already vanish in
$R_k$. Therefore the tensor expression was zero. Reversing the invertible
denominator multiplications proves injectivity.

The fraction field of $R_E$ is a subfield of $\mathscr M(B)$ containing
$\C$, so its common constants are exactly $\C$. It is generated by a
fundamental matrix, hence is Picard--Vessiot by \cref{input:pv}.
For the group identification, the Picard--Vessiot torsor identity for
$R_k$ is
\[
 R_k\otimes_K R_k\simeq R_k\otimes_k k[H_v].
\]
Tensor it with $E$ and use \eqref{eq:pv-extension}. This is the torsor
identity for $R_E$ with group $(H_v)_\C$. The action sends $Y$ to $Yh$.
Equivalently, the matrix relation ideal and its tensor symmetries have
just been shown to commute with this constant extension.
\end{proof}

\begin{remark}[Why this base-change argument is needed]
An exact faithful functor on the original $k$-category does not by itself
identify every later complex differential subobject. Here that step is
proved using a normalized solution with algebraic Taylor coefficients.
Actual period matrices need not have algebraic Taylor coefficients; they
are introduced only after \cref{prop:pv-base-change}.
\end{remark}

\section{The torsor and the exact constant-field argument}\label{sec:torsor}

We now end with the complete comparison and descent argument. The
mathematics of this section is also an abstract implication: it applies to
any finite-type neutral Tannakian system with the fiber functors and the
exactness established in \cref{sec:galois}, and with the algebraic
connection realization used in \cref{sec:pv-base-change}.

\subsection{The actual constant-period base}

Let $\omega_{\dR,0}:\Tcal_0\to\VecCat_k$ be the constant de Rham fiber
of \cref{cor:constants}. Define the comparison torsors
\begin{align*}
 X&=\Isom^\otimes(\omega_\B\otimes_k K,\omega_\dR),\\
 X_0&=\Isom^\otimes(\omega_\B|_{\Tcal_0},\omega_{\dR,0}).
\end{align*}
Thus $X$ is over $K$, whereas $X_0$ is over $k$. In a fiber product the
base change $X_{0,K}$ is always understood explicitly. These are affine
torsors by \cite[Theorem~3.2]{DeligneMilne}.

Let $p_0\in X_0(\C)$ be the actual constant comparison point and put
\begin{equation}\label{eq:actual-base}
 \begin{split}
 Z_0&=\overline{\{p_0\}}^{\mathrm{Zar}},\qquad
 A_0=k[Z_0]=k[X_0]/\ker(\operatorname{ev}_{p_0}),\\
 F_0&=\Frac(A_0),\qquad R_0=K\otimes_k A_0,\qquad
 K_0=\Frac(R_0)=F_0(S).
 \end{split}
\end{equation}
Both $A_0$ and $F_0$ are identified with their actual embedded values in
$\C$. This imposes their relations without asserting that the unknown
constant-period relations are absent. All derivations kill $F_0$.

\begin{lemma}[Domains, embeddings, and base constants]\label{lem:base-field}
There are injective maps
\[
 R_0\hookrightarrow K_0\hookrightarrow E=\C(S),
\]
and the common constant field of $K_0$ is exactly $F_0$.
\end{lemma}
\begin{proof}
Evaluation makes $A_0$ a domain. The smooth connected variety $S$ is
geometrically integral over the algebraically closed field $k$. After any
field extension its function-field construction is therefore a domain.
In particular $K\otimes_k\C$ is a domain whose fraction field is $E$.
Flatness of $K$ over $k$ embeds $R_0$ in $K\otimes_k\C$, giving the
claimed injections and identifying $K_0$ with $F_0(S)$.

The derivations $\partial_i$ form a basis of the tangent derivations of
$K_0/F_0$. If they all kill $u\in K_0$, then $du=0$ in
$\Omega^1_{K_0/F_0}$. In characteristic zero this implies that $u$ is
algebraic over $F_0$. Geometric integrality of $S_{F_0}$ implies that
$F_0$ is relatively algebraically closed in $F_0(S)$, so $u\in F_0$.
The reverse inclusion is part of the definition of the derivations.
\end{proof}

Set
\begin{equation}\label{eq:relative-scheme}
 X_{\mathrm{rel}}=X\times_{X_{0,K}}(\Spec K\times_k Z_0),
 \qquad B_0=\Ocal(X_{\mathrm{rel}}).
\end{equation}

\begin{proposition}[The correct group and its torsor]\label{prop:relative-torsor}
The map $X_{\mathrm{rel}}\to\Spec R_0$ is a torsor under
$H_{\B,R_0}$. Its coordinate algebra $B_0$ is faithfully flat and
torsion-free over $R_0$.
\end{proposition}
\begin{proof}
The normal subgroup $H_\B$ of \eqref{eq:betti-groups} is defined over
$k$. The map $X\to X_{0,K}$ is an $H_{\B,K}$-torsor: after a faithfully
flat trivialization of $X$, it is the quotient map $G_\B\to G_{0,\B}$,
whose kernel is $H_\B$. Descend this assertion and base change along
$\Spec K\times_k Z_0\to X_{0,K}$. An affine torsor under an affine
algebraic group is faithfully flat. Since $R_0$ is a domain, flatness
implies torsion-freeness.
\end{proof}

There is a natural differential structure on $B_0$. For a tensor-generator
connection in the frame of \eqref{eq:system}, its universal comparison
matrix $U$ satisfies
\begin{equation}\label{eq:universal-derivation}
 \partial_iU=A_iU.
\end{equation}
Naturality and tensor relations are stable under these operators because
the de Rham morphisms and tensor structure are horizontal. Constant
comparison coefficients are horizontal; hence restriction along $Z_0$
also respects the derivations. Thus actual analytic comparison defines a
differential map $B_0\to\mathscr M(B)$ without yet claiming it injective.

\subsection{An explicit equivariant comparison over \texorpdfstring{$\C(S)$}{C(S)}}

Use the tensor generator, algebraic frame, and normalized matrix $Y$ of
\cref{sec:pv-base-change}. Let $P$ be its actual comparison matrix from
Betti to de Rham coordinates, and let $c=P(s_0)$ be the matrix of the
comparison in \cref{lem:conjugation}. Uniqueness of solutions gives
\begin{equation}\label{eq:P-Yc}
 P=Yc.
\end{equation}
The coefficients of $c$ belong to $\C$ and need not belong to $F_0$.
This fact is allowed; those coefficients are not adjoined to the smaller
solution field in the descent argument.

\begin{lemma}[The analytic algebra and equivariance]\label{lem:analytic-torsor}
Evaluation gives a surjective $E$-algebra map
\begin{equation}\label{eq:evaluation-E}
 B_0\otimes_{R_0}E\longrightarrow E[P,\det(P)^{-1}]=R_E.
\end{equation}
The induced closed immersion of schemes is equivariant for
$H_{\B,E}$, under the identification
$(H_\B)_\C=c^{-1}(H_v)_\C c$.
\end{lemma}
\begin{proof}
The actual comparison is natural and tensor-compatible. Its restriction to
constant objects is $p_0$, so it factors through $X_{\mathrm{rel}}$.
Because $\Tcal$ has a tensor generator, its comparison algebra is generated
by the entries of the generator's universal matrix and its inverse
determinant, with its tensor relations. Evaluation therefore has precisely
the image on the right of \eqref{eq:evaluation-E}, and is surjective.
By \eqref{eq:P-Yc}, $E[P,\det(P)^{-1}]=E[Y,\det(Y)^{-1}]$.

The group of the latter Picard--Vessiot algebra acts by $Y\mapsto Yh$,
$h\in(H_v)_\C$. Consequently its action on $P$ is
\[
 P=Yc\longmapsto Yhc=P(c^{-1}hc).
\]
By \cref{lem:conjugation}, $c^{-1}hc$ ranges over $(H_\B)_\C$.
On the motivic comparison torsor, precomposition by a Betti tensor
automorphism is exactly right multiplication $U\mapsto Ug$.
The two formulas coincide. They are identities for the universal matrix
coordinate functions, so they establish equivariance as group-scheme
actions, not only on a selected set of analytic points.
\end{proof}

\begin{lemma}[Equivariant torsor maps]\label{lem:torsor-iso}
An equivariant morphism between two nonempty torsors under the same affine
group over a field is an isomorphism.
\end{lemma}
\begin{proof}
Make a faithfully flat extension that trivializes both torsors and choose
origins. An equivariant map between the two right regular group torsors
is left translation by the image of the identity, hence is invertible.
The inverse descends, since invertibility of a morphism is faithfully
flat local on the base.
\end{proof}

\begin{theorem}[Injective relative period comparison]\label{thm:period-injective}
The actual map $B_0\to\mathscr M(B)$ is injective. Moreover
\begin{equation}\label{eq:complex-isomorphism}
 B_0\otimes_{R_0}\C(S)\simeq R_E
\end{equation}
as differential algebras and equivariant torsor coordinate algebras.
\end{theorem}
\begin{proof}
By \cref{prop:pv-base-change}, the source scheme of the closed immersion
in \cref{lem:analytic-torsor} is a Picard--Vessiot torsor under
$H_{\B,E}$, in the Betti coordinates just computed. By
\cref{prop:relative-torsor}, its target is a torsor under the same group.
\Cref{lem:torsor-iso} makes that immersion an isomorphism. This proves
\eqref{eq:complex-isomorphism}.

By \cref{prop:relative-torsor} and \cref{lem:base-field}, the actual
period map is the composite of injections
\[
 B_0\ \lhook\joinrel\longrightarrow\ B_0\otimes_{R_0}K_0
 \ \lhook\joinrel\longrightarrow\ B_0\otimes_{R_0}E
 \ \xrightarrow{\sim}\ R_E
 \ \lhook\joinrel\longrightarrow\ \mathscr M(B).
\]
The first injection is localization of a torsion-free module over a
domain. The second is scalar extension along the field embedding
$K_0\hookrightarrow E$, hence is injective. The last is the concrete
analytic Picard--Vessiot embedding. This proves injectivity without
assuming anything about the constant field of the smaller period-generated
solution field.
\end{proof}

\subsection{Descent of differential simplicity and constants}

Put
\[
 R=B_0\otimes_{R_0}K_0,\qquad L=\Frac(R).
\]
These are domains by \cref{thm:period-injective}. In particular
$R=K_0[P,\det(P)^{-1}]$ in the actual analytic embedding, not an algebra
obtained by adjoining the separate entries of $c$.

\begin{lemma}[Intersection in a faithfully embedded tensor product]\label{lem:intersection}
Let $F\subset L_1,L_2$ be fields. In $L_1\otimes_F L_2$, equality
$x\otimes1=1\otimes y$ with $x\in L_1$ and $y\in L_2$ implies
$x=y\in F$. The same conclusion holds if the equality is checked in any
field into which this tensor product injects.
\end{lemma}
\begin{proof}
If $x\notin F$, the vectors $1,x$ are linearly independent over $F$.
Choose an $F$-linear functional $\ell:L_1\to F$ with
$\ell(1)=0$ and $\ell(x)=1$, extending from their span to a basis.
Applying $\ell\otimes\id$ to the equality gives $1=0$, a contradiction.
Thus $x\in F$, and the equality then implies $y=x$ by injectivity of the
field into its scalar extension. The last assertion reduces to this one
by the assumed injection.
\end{proof}

\begin{theorem}[Picard--Vessiot descent over the actual constant field]\label{thm:pv-descent}
The algebra $R$ is differentially simple, its fraction field satisfies
\begin{equation}\label{eq:exact-constants}
 L^{\partial_1=\cdots=\partial_d=0}=F_0,
\end{equation}
and $R$ is the Picard--Vessiot algebra for the realized tensor generator
over $K_0=F_0(S)$, with group $(H_\B)_{F_0}$.
\end{theorem}
\begin{proof}
By \cref{thm:period-injective},
$R\otimes_{K_0}E\simeq R_E$ is differentially simple. If $I\subset R$
is a nonzero differential ideal, faithful flatness of the field extension
makes $I\otimes E$ nonzero. Simplicity makes it all of $R_E$.
Applying faithful flatness to $(R/I)\otimes E=0$ gives $I=R$.
Thus $R$ is differentially simple.

To prove \eqref{eq:exact-constants}, localize the isomorphism
$R\otimes_{K_0}E\simeq R_E$ at every nonzero element of $R$. Since
all those elements map to nonzero elements of the domain $R_E$, this gives
an injection
\begin{equation}\label{eq:localized-tensor}
 L\otimes_{K_0}E\lhook\joinrel\longrightarrow\Frac(R_E)
 \lhook\joinrel\longrightarrow\mathscr M(B).
\end{equation}
Let $b\in L$ be killed by every $\partial_i$. Its analytic image is a
constant $\alpha\in\C\subset E$: on a connected polydisc a meromorphic
function with all partial derivatives zero is constant. In the last field
of \eqref{eq:localized-tensor} the two elements $b\otimes1$ and
$1\otimes\alpha$ have equal image. The injection gives their equality
in the tensor product. \Cref{lem:intersection} implies $b\in K_0$.
Its derivatives are zero there, so \cref{lem:base-field} gives
$b\in F_0$. Conversely, $F_0$ is killed by all the derivations.

Finally $R$ is generated over $K_0$ by the actual fundamental matrix
$P$ and its inverse determinant, it is differentially simple, and its
fraction field has exactly the constants $F_0$. These are the
Picard--Vessiot properties over this possibly non-algebraically-closed
constant field. The already established torsor coaction is defined over
$F_0$ and identifies its group with $(H_\B)_{F_0}$; equivalently this
identification follows after the faithfully flat constant extension to
$\C$ from \cref{lem:analytic-torsor} and descends. This completes the
constant-field argument rather than presupposing it.
\end{proof}

For $d=0$, $S=\Spec k$, the relative group is trivial and the argument
reduces to the embedding $A_0\hookrightarrow\C$; set $K_0=L=F_0$.
For $d>0$, the proofs above use every base derivation. Injectivity on $B$
implies injectivity on any common analytic continuation domain, so the
same result holds on the universal-cover formulation of the period map.
Every finite packet in the generic envelope lies in a category of the form
used here. Thus these replacement arguments give the intended generic
functional conclusion, using the standard external inputs
and the realization interface constructed in
\cref{thm:algebraic-interface,thm:algebraic-exponential,prop:relative-Betti-interface}.
They give neither a new algebraic-fiber
specialization theorem nor an elementary presentation of all motivic
relations. A fresh audit must check the new assembly proofs and the complete
sequence of implications before this note is integrated into Paper IV.

\section*{Research disclosure}
The authors used AI-assisted tools for proof development, source checking,
and document preparation. The authors are responsible for the note. The
frozen baseline and the repair ledger record the distinction between a
proposed replacement proof, a build check, and an independent mathematical
audit.

\small
\begin{thebibliography}{99}
\bibitem{Andre}
Y.~Andr\'e, \emph{Mumford--Tate groups of mixed Hodge structures and the
theorem of the fixed part}, Compositio Math. \textbf{82} (1992), 1--24.
\bibitem{AyoubBarbieriViale}
J.~Ayoub and L.~Barbieri-Viale, \emph{Nori $1$-motives}, Math. Ann.
\textbf{361} (2015), 367--402. \href{https://arxiv.org/abs/1206.5923}{arXiv:1206.5923}.
\bibitem{BarbieriVialeSrinivas}
L.~Barbieri-Viale and V.~Srinivas, \emph{Albanese and Picard $1$-motives},
M\'em. Soc. Math. Fr. (N.S.) \textbf{87} (2001), vi+104 pp.
\bibitem{CHM}
T.~Crespo, Z.~Hajto, and R.~Mohseni, \emph{Real Liouvillian extensions of
partial differential fields}, SIGMA \textbf{17} (2021), 095, especially
\S2.1 and its comparison with Kolchin's integrable-system theory.
\href{https://doi.org/10.3842/SIGMA.2021.095}{doi:10.3842/SIGMA.2021.095}.
\bibitem{DeligneMilne}
P.~Deligne and J.~S.~Milne, \emph{Tannakian categories}, in
\emph{Hodge Cycles, Motives, and Shimura Varieties}, Lecture Notes in Math.
\textbf{900}, Springer, 1982, 101--228. Author-hosted corrected text:
\url{https://www.jmilne.org/math/xnotes/tc.pdf}.
\bibitem{DouaiSabbah}
A.~Douai and C.~Sabbah, \emph{Gauss--Manin systems, Brieskorn lattices
and Frobenius structures (I)}, Ann. Inst. Fourier \textbf{53} (2003),
1055--1116, Propositions~1.18 and~1.20.
\bibitem{GallauerPepin}
M.~Gallauer and S.~Pepin Lehalleur, \emph{Exponentiation of coefficient
systems and exponential motives}, \href{https://arxiv.org/abs/2211.17247}{arXiv:2211.17247}.
\bibitem{HuberWustholz}
A.~Huber and G.~W\"ustholz, \emph{Transcendence and Linear Relations of
$1$-Periods}, Cambridge Tracts in Mathematics \textbf{227}, Cambridge
University Press, 2022. \href{https://arxiv.org/abs/1805.10104}{arXiv:1805.10104},
Theorem~9.14.
\bibitem{IvorraMorel}
F.~Ivorra and S.~Morel, \emph{The four operations on perverse motives},
J. Eur. Math. Soc. \textbf{26} (2024), 4191--4272.
\bibitem{JacobsenTerenzi}
E.~Jacobsen and L.~Terenzi, \emph{A comparison of categories of Nori
motivic sheaves}, \href{https://arxiv.org/abs/2509.21476}{arXiv:2509.21476},
Corollary~5.15.
\bibitem{Kolchin}
E.~R.~Kolchin, \emph{Picard--Vessiot theory of partial differential fields},
Proc. Amer. Math. Soc. \textbf{3} (1952), 596--603.
\href{https://doi.org/10.1090/S0002-9939-1952-0049883-8}{doi:10.1090/S0002-9939-1952-0049883-8}.
\bibitem{LamLitt}
Y.-H.~J.~Lam and D.~Litt, \emph{Geometric local systems on the projective
line minus four points}, Compos. Math. \textbf{161} (2025), 536--554,
Lemma~3.5. \href{https://doi.org/10.1112/S0010437X24007620}{doi:10.1112/S0010437X24007620}.
\bibitem{Mostaed}
A.~Mostaed, \emph{On the invariant cycle theorem for families of Nori
motives}, \href{https://arxiv.org/abs/2206.05582}{arXiv:2206.05582}, \S4.
\bibitem{PutSinger}
M.~van der Put and M.~F.~Singer, \emph{Galois Theory of Linear Differential
Equations}, Grundlehren Math. Wiss. \textbf{328}, Springer, 2003, Chapter~1.
\href{https://doi.org/10.1007/978-3-642-55750-7}{doi:10.1007/978-3-642-55750-7}.
\bibitem{SabbahMonodromy}
C.~Sabbah, \emph{Monodromy at infinity and Fourier transform II},
Publ. Res. Inst. Math. Sci. \textbf{42} (2006), 803--835.
Proposition~5.8, Corollary~5.20, and (5.21)--(5.22). Author-hosted text:
\url{https://perso.pages.math.cnrs.fr/users/claude.sabbah/articles/sabbah_Fourier-twistor.pdf}.
\bibitem{SabbahParameters}
C.~Sabbah, \emph{A short proof of a theorem of Cotti, Dubrovin and
Guzzetti}, Port. Math. \textbf{79} (2022), 29--43.
\href{https://arxiv.org/abs/2103.16878v3}{arXiv:2103.16878v3},
Appendix~B, preprint pp.~12--13 (the assertion immediately before
Proposition~B.1 and that proposition).
\bibitem{Snodgrass}
B.~Snodgrass, \emph{Periods of $E$-operators},
\href{https://arxiv.org/abs/2608.06005}{arXiv:2608.06005} (2026),
Corollary~2.17 and Theorem~3.6.
\bibitem{Terenzi}
L.~Terenzi, \emph{Tensor structure on perverse Nori motives},
Ann. K-Theory \textbf{11} (2026), 47--170.
\href{https://arxiv.org/abs/2401.13547}{arXiv:2401.13547}.
\bibitem{Tubach}
S.~Tubach, \emph{On the Nori and Hodge realisations of Voevodsky motives},
Compos. Math. \textbf{161} (2025), 2155--2201.
\href{https://arxiv.org/abs/2309.11999}{arXiv:2309.11999}.
\bibitem{Virk}
R.~Virk, \emph{Realizations of exponential sheaves and Fourier transform},
\href{https://arxiv.org/abs/2605.14904v4}{arXiv:2605.14904v4} (2026),
\S\S2.6--2.7 and~6.6--6.7, Propositions~6.4, 6.8, and~7.2,
Theorems~6.13 and~7.5--7.6, and Appendix~B.

\bibitem{ChoudhuryGallauer}
U.~Choudhury and M.~Gallauer Alves de Souza,
\emph{An isomorphism of motivic Galois groups},
Adv. Math. \textbf{313} (2017), 470--536.
\href{https://arxiv.org/abs/1410.6104}{arXiv:1410.6104},
\S\S6--7, Propositions~7.1 and~7.12.
\bibitem{DrewGallauer}
B.~Drew and M.~Gallauer,
\emph{The universal six-functor formalism},
\href{https://arxiv.org/abs/2009.13610v4}{arXiv:2009.13610v4},
Definitions~7.5 and~7.7, Theorem~7.14.
\bibitem{Harrer}
D.~Harrer,
\emph{Comparison of the categories of motives defined by Voevodsky and Nori},
Ph.D. thesis, University of Freiburg, 2016.
\href{https://arxiv.org/abs/1609.05516}{arXiv:1609.05516},
Chapters~4 and~7; Theorems~7.4.17 and~7.6.5.
\bibitem{SaitoMixed}
M.~Saito, \emph{On the formalism of mixed sheaves},
Preprint RIMS-784, August 1991; electronic version 2006.
\href{https://arxiv.org/abs/math/0611597}{arXiv:math/0611597},
\S1.8(ii), \S\S2--6, and \S6.15.
\bibitem{FresanJossen}
J.~Fres\'an and P.~Jossen, \emph{Exponential motives},
preliminary manuscript, 2020 version, with the absolute rapid-decay
comparison used here also recalled in \cite[\S3]{Snodgrass}.
\bibitem{Hien}
M.~Hien, \emph{Periods for flat algebraic connections},
Invent. Math. \textbf{178} (2009), 1--22.
\bibitem{HienRoucairol}
M.~Hien and C.~Roucairol,
\emph{Integral representations for solutions of exponential Gauss--Manin systems},
Bull. Soc. Math. France \textbf{136} (2008), 505--532.

\end{thebibliography}
\end{document}
